%% kapa: exam=1 round=28 year=2017 period=2 form=A subject=appraisal_law subject_ko=감정평가관계법규 seq=31 number=31
%% source: Q-Net 감정평가사 기출문제 (cst00309) artlSeq=5199603 file=extracted/1008/5199603/2017년 제28회 감정평가사 1차시험 문제지/2교시 A형.pdf
%% conversion: kordoc 4.24.51 -> tools/gen1.py (auto); review=proofread-sonnet; flags=none
\begin{question}{31}
공간정보의 구축 및 관리 등에 관한 법령상 토지의 등록에 관한 설명으로 옳지 않은 것은?
\begin{choices}
  \item 국토교통부장관은 모든 토지에 대하여 필지별로 소재ㆍ지번ㆍ지목ㆍ면적ㆍ경계 또는 좌표 등을 조사ㆍ측량하여 지적공부에 등록하여야 한다.
  \item 지번은 지적소관청이 지번부여지역별로 차례대로 부여한다.
  \item 지적공부에 등록하는 경계 또는 좌표는 토지의 이동이 있을 때 토지소유자의 신청이 없는 경우 지적소관청이 직권으로 조사ㆍ측량하여 결정할 수는 없다.
  \item 물을 상시적으로 이용하지 않고 약초를 주로 재배하는 토지의 지목은 전(田)이다.
  \item 지목은 필지마다 하나의 지목을 설정하여야 한다.
\end{choices}
\end{question}
\answer{3}
